\documentclass[11pt]{amsart}

\usepackage[T1]{fontenc}
\usepackage{lmodern}
\usepackage{microtype}
\usepackage{amsmath,amssymb,amsthm}
\usepackage[hidelinks]{hyperref}

\hypersetup{
  pdftitle={Disconnected Counterexamples to a Real-Rootedness Conjecture for Weighted Bond Posets}
}

\newtheorem{theorem}{Theorem}
\theoremstyle{remark}
\newtheorem{remark}[theorem]{Remark}

\title[Disconnected counterexamples]
{Disconnected Counterexamples to a Real-Rootedness Conjecture
for Weighted Bond Posets}

\date{}

\begin{document}

\begin{abstract}
Conjecture~4.13(2) of Gonz\'alez D'Le\'on and Wachs
\cite{GDLW} asserts that a signed difference of weighted-bond-poset
M\"obius polynomials is real-rooted whenever the underlying graphs
have the same numbers of vertices and connected components and one
is a subgraph of the other. We show that the conjecture is false for
disconnected graphs. For every integer $r\geq 3$, the pair
$G_r=rK_3$ and $H_r=rP_3$ is a counterexample, where $P_3$ is
the three-vertex path and $rX$ denotes the disjoint union of $r$ copies
of $X$. The connected restriction is not addressed.
\end{abstract}

\maketitle

Let $\mu_G(t)$ denote the M\"obius polynomial of the weighted bond
poset of a graph $G$, as defined in \cite{GDLW}. Conjecture~4.13(2)
of \cite{GDLW} states that if $G$ has $n$ vertices and $k$ connected
components and $H$ is a subgraph of $G$ with the same $n$ and $k$,
then
\[
(-1)^{n-k}\bigl(\mu_G(t)-\mu_H(t)\bigr)
\]
is real-rooted. The conjecture is stated for arbitrary $k$.

\begin{theorem}
For every integer $r\geq 3$, let
\[
G_r=rK_3,
\qquad
H_r=rP_3,
\]
where each copy of $P_3$ is obtained by deleting one edge from the
corresponding copy of $K_3$. Then $H_r$ is a spanning subgraph of
$G_r$, both graphs have $3r$ vertices and $r$ connected components,
and
\[
(-1)^{3r-r}\bigl(\mu_{G_r}(t)-\mu_{H_r}(t)\bigr)
\]
is not real-rooted. Consequently, Conjecture~4.13(2) of
\cite{GDLW} is false as stated.
\end{theorem}

\begin{proof}
Set
\[
A(t)=2+5t+2t^2,
\qquad
B(t)=1+3t+t^2.
\]
Equation~(1.5) and Proposition~2.2 of \cite{GDLW} give
\[
\mu_{G_r}(t)-\mu_{H_r}(t)=A(t)^r-B(t)^r.
\]
Since $3r-r=2r$, the sign in the conjecture is positive.

Let $\zeta=e^{2\pi i/r}$ and $c=\cos(2\pi/r)$. Since $r\geq3$,
$\zeta$ is nonreal and $|c|<1$. The factorization
\[
A(t)^r-B(t)^r
=\prod_{j=0}^{r-1}\bigl(A(t)-\zeta^jB(t)\bigr)
\]
shows that the real polynomial
\[
\begin{aligned}
Q_r(t)
&=\bigl(A(t)-\zeta B(t)\bigr)
  \bigl(A(t)-\overline{\zeta}B(t)\bigr)\\
&=A(t)^2-2cA(t)B(t)+B(t)^2
\end{aligned}
\]
divides $A(t)^r-B(t)^r$. For real $t$,
\[
Q_r(t)
=\bigl(A(t)-cB(t)\bigr)^2+(1-c^2)B(t)^2.
\]
Thus a real zero of $Q_r$ would be a common zero of $A$ and $B$.
But
\[
A(t)-B(t)=(1+t)^2,
\qquad
A(-1)=B(-1)=-1,
\]
so $A$ and $B$ have no common real zero. Hence $Q_r(t)>0$ for every
real $t$. Its leading coefficient is
\[
|2-\zeta|^2=5-4c>0,
\]
so $Q_r$ is a quartic polynomial with no real zeros. Since $Q_r$
divides $A^r-B^r$, the latter is not real-rooted.
\end{proof}

The first member of the family occurs at $r=3$. It has nine vertices
and three connected components, and
\[
\mu_{3K_3}(t)-\mu_{3P_3}(t)
=(1+t)^2\bigl(7+37t+63t^2+37t^3+7t^4\bigr),
\]
where the quartic factor is strictly positive on $\mathbb{R}$.

\begin{remark}
The threshold $r\geq3$ is exact within this family. Indeed,
\[
A-B=(1+t)^2,
\qquad
A^2-B^2=(1+t)^2(3+8t+3t^2),
\]
and the quadratic factor has discriminant $28$.
\end{remark}

\begin{remark}
These examples do not refute the companion $\gamma$-positivity
statement, Conjecture~4.12(2) of \cite{GDLW}. Writing
$X=(1+t)^2$, one has $A=2X+t$ and $B=X+t$, and therefore
\[
A^r-B^r
=\sum_{j=0}^{r-1}
\binom{r}{j}\bigl(2^{r-j}-1\bigr)
t^j(1+t)^{2r-2j}.
\]
Thus the difference is $\gamma$-positive for every $r\geq1$.
The construction also makes no assertion about Conjecture~4.13(1)
or the restriction of Conjecture~4.13(2) to connected graphs.
\end{remark}

\begin{thebibliography}{1}

\bibitem{GDLW}
R.~S. Gonz\'alez D'Le\'on and M.~L. Wachs,
\emph{Weighted bond posets and a new chromatic symmetric function},
\href{https://arxiv.org/abs/2608.08692v1}
{arXiv:2608.08692v1 [math.CO]} (2026).

\end{thebibliography}

\end{document}